
\documentclass{article}
\usepackage{colortbl}
\usepackage{makecell}
\usepackage{multirow}
\usepackage{supertabular}

\begin{document}

\newcounter{utterance}

\centering \large Interaction Transcript for game `dond', experiment `coop\_en', episode 0 with Qwen3{-}32B{-}t0.0.
\vspace{24pt}

{ \footnotesize  \setcounter{utterance}{1}
\setlength{\tabcolsep}{0pt}
\begin{supertabular}{c@{$\;$}|p{.15\linewidth}@{}p{.15\linewidth}p{.15\linewidth}p{.15\linewidth}p{.15\linewidth}p{.15\linewidth}}
    \# & $\;$A & \multicolumn{4}{c}{Game Master} & $\;\:$B\\
    \hline

    \theutterance \stepcounter{utterance}  
    & & \multicolumn{4}{p{0.6\linewidth}}{
        \cellcolor[rgb]{0.9,0.9,0.9}{
            \makecell[{{p{\linewidth}}}]{
                \texttt{\tiny{[P1$\langle$GM]}}
                \texttt{You are playing a cooperative negotiation game in which you have to agree on how to divide a set of items among you and another player.} \\
\\ 
\texttt{Rules:} \\
\texttt{(a) You and the other player are given a set of items. Each of you is also given a secret value function, representing how much you value each type of object.} \\
\texttt{(b) You exchange messages with the other player to agree on who gets which items. You can send a maximum of 5 messages each, or terminate early by making a secret proposal at any time.} \\
\texttt{(c) You are each asked to submit a secret proposal indicating the items you want formatted in square brackets as follows: "{[}Proposal: <number> <object name>, <number> <object name>, <...>{]}"} \\
\texttt{(d) If your proposals are complementary, i.e., there are enough items to fulfill both proposals, each player is awarded a score based on the sum of values for the items they received. Otherwise, both of you get zero points.} \\
\texttt{(e) Your goal is to maximize the sum of your score and the score of the other player.} \\
\\ 
\texttt{Let us start.} \\
\\ 
\texttt{The set of available items is:} \\
\\ 
\texttt{3 rings, 1 dice, 1 pan} \\
\\ 
\texttt{Your secret value function is:} \\
\\ 
\texttt{ring: 0, dice: 4, pan: 6} \\
\\ 
\texttt{IMPORTANT: Your messages, unless it is the secret proposal, are directly transmitted to the other player, so do not include any response to the rules or text announcing your message. To make a secret proposal, use the indicated format. Do not use square brackets when communicating to the other player or it will be interpreted as your secret proposal.} \\
            }
        }
    }
    & \\ \\

    \theutterance \stepcounter{utterance}  
    & \multicolumn{4}{p{0.6\linewidth}}{
        \cellcolor[rgb]{0.9,0.9,0.9}{
            \makecell[{{p{\linewidth}}}]{
                \texttt{\tiny{[P1$\rangle$GM]}}
                \texttt{<think>} \\
\texttt{Okay, let's see. I need to divide the items between me and another player. The items are 3 rings, 1 dice, and 1 pan. My value function is ring: 0, dice: 4, pan: 6. So I don't value rings at all, but the dice and pan are worth more. The other player's value function is secret, so I don't know what they value. The goal is to maximize the total score for both of us.} \\
\\ 
\texttt{First, I should consider what the other player might want. Since there are 3 rings, maybe they value rings a lot. If they do, then I should let them take the rings, and I can take the dice and pan. But wait, there's only one dice and one pan. So if I take both, that's possible, but I need to check if the other player's proposal is complementary.} \\
\\ 
\texttt{But how do I communicate this? I can send up to 5 messages. Let me think. Maybe I should propose a division where I take the dice and pan, and they take all the rings. That way, I get 4 + 6 = 10 points, and if they value rings, they might get points from the rings. The total would be 10 plus their points. But since I don't know their values, maybe I should ask them what they prefer. But I can't ask directly because the messages are exchanged. Wait, the messages are directly transmitted, so I can ask them to indicate their preferences. But they might not respond honestly. Hmm.} \\
\\ 
\texttt{Alternatively, maybe I can make a secret proposal. Let me think. If I propose taking the dice and pan, and the other player proposes taking the rings, then our proposals are complementary. Since there are 3 rings, if they want all 3, that's possible. But if they also want the dice or pan, there's a conflict. But since I don't know their values, maybe the best approach is to propose taking the items I value (dice and pan) and let them take the rings. That way, if they value rings, they can take them, and we both get our desired items. } \\
\\ 
\texttt{Wait, but the secret proposal has to be in the format {[}Proposal: ...{]}. So I need to make sure that my proposal is such that if the other player's proposal doesn't conflict, we both get points. Let me check} \\
            }
        }
    }
    & & \\ \\

    \theutterance \stepcounter{utterance}  
    & & & \multicolumn{2}{p{0.3\linewidth}}{
        \cellcolor[rgb]{0.9,0.9,0.9}{
            \makecell[{{p{\linewidth}}}]{
                \texttt{\tiny{[GM$|$GM]}}
                \texttt{proposal does not include number/name pairs} \\
            }
        }
    }
    & & \\ \\

\end{supertabular}
}

\end{document}
